% Source pp. 182-200 (Lectures13-18.pdf).
\section{Fixed Points of RG}

At generic point in $\Sigma$ the pattern of the RG flow is rather boring: it just
flows. Something interesting is expected to happen near singular points of the
flow. Of these the simplest and probably the most important ones are so called
\emph{Fixed Points}.

A fixed point is an action $\mathcal A_*$ which satisfies the equation
\[
B\{\mathcal A_*\}=0\,.
\]
Associated trajectory is trivial,
\[
\mathcal A_l=\mathcal A_*\qquad\text{for all}\quad l\,.
\]
As the correlation length measured in the units $\Lambda$ is
\[
R_c(l)=e^{-l}\,R_c
\]
for $\mathcal A_l$, this means that for $\mathcal A_*$ we must have
\[
\text{either}\quad R_c=0\qquad\text{or}\quad R_c=\infty\,.
\]
By the definition, $\mathcal A_*$ with $R_c=\infty$ belongs to the ``critical
surface'',
\[
\mathcal A_*\in\Sigma_c\qquad\text{if}\qquad R_c=\infty\,.
\]
We call the fixed points $\mathcal A_*$ with $R_c=\infty$ the ``critical'' fixed
points.

To illustrate the topological ideas related to the Wilson's RG, I will describe
the simplest possible form of topology. There are many ways in which topology of
the RG flow can be more complicated but complications can be understood once the
simplest picture is grasped. In this analysis I will assume continuity of the RG
flow, i.e.\ the functions $B$ in \eqref{eq:Bfunctions} are continuous functions of
their arguments. Again, general proof of this statement is not available, but
various examples support this assumption.

Assume there are only two fixed points, one with $R_c=0$ and one with
$R_c=\infty$. Denote them $P_0$ and $P_\infty$ (figure~\ref{fig:twofixed}, left).
\begin{figure}[ht]
\centering\figfile{twofixed}
\caption{Left: surfaces of constant $R_c$ with the two fixed points $P_0$
($R_c=0$) and $P_\infty$ on the critical surface ($R_c=\infty$). Right: the
trajectory $U$ leaving $\Sigma_c$ at $P_\infty$, and a generic trajectory $G$
starting close to $\Sigma_c$.}
\label{fig:twofixed}
\end{figure}

As $R_c(l)=e^{-l}R_c$, any trajectory which starts at finite $R_c$ will approach
$R_c=0$ and then $P_0$. On the other hand trajectories which start on $\Sigma_c$
will approach $P_\infty$. The only point where some trajectory can leave
$\Sigma_c$ is $P_\infty$.\footnote{Suppose a trajectory ``leaves'' $\Sigma_c$ at
some $P_c\in\Sigma_c$. Obviously $P_c=\lim_{l\to\infty}RG_{-l}\{P\}$ for any $P$ on
such trajectory.} So there must be a trajectory (let's call it $U$) going from
$P_\infty$ to $P_0$ (figure~\ref{fig:twofixed}, right).

Consider also a trajectory $G$ which starts close to $\Sigma_c$. This one first
approaches $P_\infty$ while staying close to $\Sigma_c$, then spends very long
``time'' $l$ in the vicinity of $P_\infty$ (by continuity $B\{\mathcal A\}$ is
small when $\mathcal A$ is close to $P_\infty$) and then departs along some way
close to $U$ (recall that $U$ is the only trajectory that leaves $\Sigma_c$),
finally going towards $P_0$.

Now let us look for the space $\Sigma(\infty)$ in this picture. It is clear that
$\Sigma(\infty)$ is $U$, with the points $P_\infty$ and $P_0$ included. Starting
with any action $\mathcal A$ on the curve $U$ one can go back along the trajectory
arbitrary amount of ``time'' $l$ without leaving the space $\Sigma$ (actually,
without leaving $U$). The reason is that $R_c=\infty$ on $U$ only at the point
$P_\infty$. To reach $P_\infty$ from any point on $U$ with finite $R_c$ takes
infinite ``time'' $l$ (when moving backwards). Starting from other point in
$\Sigma$ and integrating backwards, one generally leaves $\Sigma$ in sufficient
amount of ``time''.

We see from this example the role played by critical fixed points in the
formation of the space $\Sigma(\infty)$. Therefore finding critical fixed points
is of central importance in the RG approach.

\subsection*{Critical Fixed Points}

Suppose $\mathcal A_*$ is a critical fixed point, i.e.
\[
B\{\mathcal A_*\}=0\,.
\]
and associated $R_c=\infty$. Consider vicinity of $\mathcal A_*$. Let
\[
\mathcal A=\mathcal A_*+\delta\mathcal A\,,
\]
where $\delta\mathcal A$ is infinitesimal. The variation $\delta\mathcal A$ must be
thought of as a vector in the tangent space $T\Sigma$ taken at the point
$\mathcal A_*$ (figure~\ref{fig:tangent}),
\[
\delta\mathcal A\in T\Sigma|_{\mathcal A_*}\,.
\]
\begin{figure}[ht]
\centering\figfile{tangent}
\caption{A variation $\delta\mathcal A$ of the fixed-point action is a vector in
the tangent space $T\Sigma|_{\mathcal A_*}$.}
\label{fig:tangent}
\end{figure}
One can linearize the RG transformation near $\mathcal A_*$,
\[
B\{\mathcal A_*+\delta\mathcal A\}=K\{\delta\mathcal A\}+O(\delta\mathcal A^2)\,,
\]
where $K$ is a linear operator acting in this tangent space,
\[
K:\ T\Sigma|_{\mathcal A_*}\to T\Sigma|_{\mathcal A_*}\,.
\]
As it is conventional in the theory of linear operators, I will write
\begin{equation}\label{eq:Klin}
K\{\delta\mathcal A\}\equiv K\,\delta\mathcal A\,.
\end{equation}
The RG flow in this close vicinity of $\mathcal A_*$ is then described by linear
differential equation
\[
\frac d{dl}\,(\delta\mathcal A)=K\,\delta\mathcal A\,.
\]
This equation is easily solved provided one can find eigenvectors $\Psi_\alpha$
and associated eigenvalues $\kappa_\alpha$ of the operator $K$,
\[
K\,\Psi_\alpha=\kappa_\alpha\,\Psi_\alpha\,.
\]
The solution takes the form
\begin{equation}\label{eq:linsol}
\delta\mathcal A(l)=\sum_\alpha C_\alpha\,e^{\kappa_\alpha l}\,\Psi_\alpha\,,
\end{equation}
where the constants $C_\alpha$ are determined by the initial vector
$\delta\mathcal A(0)=\delta\mathcal A_0$.

It is easy to understand that the pattern of the RG flow in the vicinity of
$\mathcal A_*$ is directly related to the eigenvalues $\kappa_\alpha$. To
illustrate this let us assume that there are only two eigenvalues (this
corresponds to a two-dimensional approximation to $\Sigma$), both real, one
positive and one negative,
\[
\kappa_+>0\,,\qquad\kappa_-<0\,.
\]
Figure~\ref{fig:saddle} shows the pattern of corresponding RG flow in the vicinity
of the fixed point. There is a natural coordinate system on
$T\Sigma|_{\mathcal A_*}$ associated with the eigenvectors $\Psi_+$ and $\Psi_-$.
\begin{figure}[ht]
\centering\figfile{saddle}
\caption{RG flow near a critical fixed point with one positive and one negative
eigenvalue. The $\Psi_-$ axis lies in the critical surface $\Sigma_c$ and flows
into the fixed point; the $\Psi_+$ axis is the trajectory $U$ leaving it. A
generic trajectory $G$ passes near the fixed point and leaves along $U$.}
\label{fig:saddle}
\end{figure}

All trajectories starting at the $\Psi_-$ axis never leave it; they flow towards
the fixed point $\mathcal A_*$. As the correlation length can only decrease with
$l$, the entire axis $\Psi_-$ must belong to $\Sigma_c$. The trajectories
starting at the $\Psi_+$ also remain at this axis, but they flow away from the
fixed point. For sufficiently large $l$ these trajectories leave the domain of
validity of the linear approximation \eqref{eq:Klin}, and \eqref{eq:linsol} says
nothing about their ultimate destination. The same is true for any generic
trajectory $G$ starting at neither axis: at sufficiently large $l$ it leaves the
vicinity of the fixed point and its fate cannot be predicted within the linear
approximation. However there is an important difference between the generic
trajectories like $G$ and any trajectory $U$ that lies on the axis $\Psi_+$.
Starting at any point on the axis $\Psi_+$ one can ``integrate backward'' for
indefinite negative $l$ without leaving the vicinity of the fixed point
$\mathcal A_*$; in fact corresponding $\mathcal A_{-l}$ approaches $\mathcal A_*$
as $l\to\infty$. This is in contrast with the negative $l$ behavior of generic
trajectories; associated $\mathcal A_{-l}$ leave the vicinity of $\mathcal A_*$
for sufficiently large $l$, and the question if they remain in $\Sigma$ for
$l\to\infty$ cannot be answered within the linear approximation. On the other
hand the linear approximation is good enough to recognize in this example at
least one-dimensional component of the subspace $\Sigma(\infty)$ associated with
the fixed point $\mathcal A_*$; the axis $\Psi_+$ is the linear approximation to
this component. Precise statement is
\[
\Psi_+\subset T\Sigma(\infty)|_{\mathcal A_*}\,.
\]

It is not difficult to generalize the conclusions obtained in the example above.
It is useful to define two special subspaces of the tangent space
$T\Sigma|_{\mathcal A_*}$. Let $\mathcal F_+$ be the linear envelope of the
eigenvectors of $K$ corresponding to positive eigenvalues $\kappa$; with some
abuse of notations we can write
\[
\mathcal F_+=\oplus_{\alpha:\kappa_\alpha>0}\,\Psi_\alpha\,.
\]
Similarly, let $\mathcal F_-$ be the sum of eigenspaces corresponding to negative
$\kappa_\alpha$, i.e.
\[
\mathcal F_-=\oplus_{\alpha:\kappa_\alpha<0}\,\Psi_\alpha\,.
\]
This separation leaves unaccounted very interesting possibility of having zero
eigenvalues, later we will discuss it separately. Also, it is assumed that all
eigenvalues of $K$ are real. At this point this assumption looks arbitrary. We
will see later that having complex eigenvalues contradicts quantum mechanical
interpretation of field theory.

Simple analysis of \eqref{eq:linsol} shows that
\begin{itemize}
\item The fixed point $\mathcal A_*$ belongs to a critical surface $\Sigma_c$ whose
dimensionality coincides with the dimensionality of $\mathcal F_-$, and
\[
\mathcal F_-=T\Sigma_c|_{\mathcal A_*}\,.
\]
\item The fixed point $\mathcal A_*$ also belongs to some component of
$\Sigma(\infty)$; the dimensionality of this component is equal or greater than
the dimensionality of $\mathcal F_+$, and
\[
\mathcal F_+\subset T\Sigma(\infty)|_{\mathcal A_*}\,.
\]
\end{itemize}

Let us define the subspace $U(\mathcal A_*)\subset\Sigma(\infty)$ by the condition
that all points of $U(\mathcal A_*)$ lie on RG trajectories which, when
``integrated backward'', converge to the fixed point $\mathcal A_*$, i.e.
\[
\mathcal A_0\in U(\mathcal A_*)\ \Rightarrow\ \mathcal A_l\to\mathcal A_*\quad\text{as}\quad l\to-\infty
\]
The subspace $U(\mathcal A_*)$ is called the \emph{unstable manifold of the fixed
point $\mathcal A_*$}. Obviously $\mathcal A_*\in U(\mathcal A_*)$. If we ignore
the possibility of zero eigenvalues then the tangent space to $U(\mathcal A_*)$
taken at the fixed point $\mathcal A_*$ itself coincides with the space
$\mathcal F_+$, i.e.
\[
\mathcal F_+=TU(\mathcal A_*)|_{\mathcal A_*}\,.
\]

The reason why $\Sigma(\infty)$ can have dimensionality greater than that of the
unstable manifold $U(\mathcal A_*)$ is illustrated by the following simple
example. Suppose we have two fixed points, the $\mathcal A_*$ and another one,
$\mathcal A'_*$, such that $\mathcal A_*\in U(\mathcal A'_*)$
(figure~\ref{fig:twofp}).
\begin{figure}[ht]
\centering\figfile{twofp}
\caption{Two fixed points with $\mathcal A_*\in U(\mathcal A'_*)$. One
trajectory of $U(\mathcal A'_*)$ flows into $\mathcal A_*$.}
\label{fig:twofp}
\end{figure}
Then there must be at least one RG trajectory in $U(\mathcal A'_*)$ which flows to
$\mathcal A_*$ (this is by the definition of unstable manifold $U$). This
trajectory has $R_c=\infty$ and therefore it simultaneously belongs to
$\Sigma(\infty)$ and $\Sigma_c$. It follows that the unstable manifold
$U(\mathcal A_*)$ has at least one dimension less than $U(\mathcal A'_*)$.

Understanding the global topological properties of the RG flow, like the one
illustrated in the above example, is beyond the capabilities of linear analysis.
However, as we have seen, it is within the power of linear analysis to prove that
each fixed point $\mathcal A_*$ generates a manifold $U(\mathcal A_*)$ which
belongs to $\Sigma(\infty)$. The dimensionality $n$ of this manifold equals the
number of positive eigenvalues of the operator $K$. In other words, the fixed
point $\mathcal A_*$ generates $n$-parameter family of finite local QFT.

The above discussion was somewhat abstract. It is useful to understand the
meaning of the abstract notions of the above linear analysis, like tangent
spaces, in terms of local field degrees of freedom.

Recall that generic quasilocal action $\mathcal A$ can be written as integral
\[
\mathcal A=\int d^dx\,\Lag(x)\,,
\]
where
\[
\Lag(x)=\Lag\bigl(\varphi_0(x),\partial_\mu\varphi_0(x),\partial_\mu\partial_\nu\varphi_0(x),\dots\bigr)
=\sum_{k=1}^\infty\frac{u_{2k}}{2k!}\,\varphi_0^{2k}(x)+\sum_{k=0}^\infty\frac{v_{2k}}{2k!}\,\varphi_0^{2k}(x)\,(\partial\varphi_0)^2(x)+\cdots
\]
is a scalar function of the $\varphi_0(x)$ and all its derivatives. The variation
$\delta\mathcal A$ has similar form,
\[
\delta\mathcal A=\int d^dx\,\delta\Lag(x)\,,\qquad
\delta\Lag(x)=\sum_{k=1}^\infty\frac{\delta u_{2k}}{2k!}\,\varphi_0^{2k}(x)+\sum_{k=0}^\infty\frac{\delta v_{2k}}{2k!}\,\varphi_0^{2k}(x)\,(\partial\varphi_0)^2(x)+\cdots\,.
\]
In other words $\delta\mathcal A$ is an integral of (infinite) linear combination
of local composite fields; we can write it as
\[
\delta\mathcal A=\int d^dx\,\sum_\alpha\delta\lambda_\alpha\,\mathcal O_\alpha(x)\,,
\]
where
$\{\mathcal O_\alpha\}=\{\varphi_0^{2n},\varphi_0^{2n}(\partial\varphi_0)^2,\dots\}$,
and $\{\delta\lambda_\alpha\}$ are variations of corresponding coupling
constants. Thus the tangent space $T\Sigma|_{\mathcal A_*}$ can be identified
with the space of scalar composite fields spanned by $\{\mathcal O_\alpha\}$. I
will use notation $\mathcal F$ for this space of local fields; we have
\[
T\Sigma|_{\mathcal A_*}=\mathcal F\,.
\]

How the RG transformation acts on these composite fields? Consider a correlation
function involving the field $\mathcal O_\alpha(x)$, say
\[
\bigl\langle\mathcal O_\alpha(x)\cdots\bigr\rangle=\frac1Z\int[D\varphi_0]\,\bigl[\mathcal O_\alpha(x)\,\varphi_0(x_1)\cdots\varphi_0(x_n)\bigr]\,e^{-\mathcal A[\varphi_0]}\,,
\]
To perform the step 1 of the Wilson's RG transformation we must split the field
$\varphi_0$ into its ``slow'' and ``fast'' parts,
\[
\varphi_0(x)=\varphi_1(x)+\tilde\varphi(x)\,,
\]
and then integrate the ``fast'' part $\tilde\varphi$ out.

Recall that $\mathcal O_\alpha(x)$ in general is a polynomial function of
$\varphi_0$ and its derivatives at the point $x$,
\[
\mathcal O_\alpha(x)=\mathcal O_\alpha\bigl(\varphi_0(x),\partial_\mu\varphi_0(x),\dots\bigr)
\]
and therefore
\[
\mathcal O_\alpha\bigl(\varphi_1(x)+\tilde\varphi(x),\ \text{derivatives}\bigr)=\sum_\beta Y_\alpha{}^\beta\bigl(\tilde\varphi(x),\ \text{derivatives}\bigr)\,\mathcal O_\beta\bigl(\varphi_1(x),\ \text{derivatives}\bigr)\,,
\]
where $Y_\alpha{}^\beta$ are polynomials of the ``fast'' field $\tilde\varphi(x)$ and
its derivatives. Example is the field $\varphi_0^n$,
\[
\bigl(\varphi_1(x)+\tilde\varphi(x)\bigr)^n=\sum_{k=0}^nY_n{}^k\bigl(\tilde\varphi(x)\bigr)\,\varphi_1^{n-k}(x)\,,\qquad Y_n{}^k(\tilde\varphi)=C_n^k\,\tilde\varphi^k\,.
\]
Assuming again that we are interested in the case when all external momenta
$p_1,\dots,p_n\ll\Lambda$, one can write the result of $\tilde\varphi$ integration
as
\begin{multline}\label{eq:Ostep1}
\frac1Z\int[D\tilde\varphi]\,\bigl[\mathcal O_\alpha\bigl(\varphi_1(x)+\tilde\varphi(x),\ \text{derivatives}\bigr)\,\varphi_0(x_1)\cdots\varphi_0(x_n)\bigr]\,e^{-\mathcal A[\varphi_1+\tilde\varphi]}\\
=\frac1{Z_1}\Bigl[\sum_\beta\bar Y_\alpha{}^\beta\bigl(L|[\varphi_1]\bigr)\,\mathcal O_\beta\bigl(\varphi_1(x),\ \text{derivatives}\bigr)\,\varphi_1(x_1)\cdots\varphi_1(x_n)\Bigr]\,e^{-\mathcal A_1[\varphi_1]}\,,
\end{multline}
where $\mathcal A_1$ is the action appearing in the step 1 of the Wilson's RG
transformation with the parameter $L$, and
\[
\bar Y_\alpha{}^\beta\bigl(L|[\varphi_1]\bigr)=\bigl\langle Y_\alpha{}^\beta(\tilde\varphi,\ \text{derivatives})\bigr\rangle_{\tilde\varphi,c}
\]
are the expectation values of $Y_\alpha{}^\beta$ over the ensemble of the fast mode
fluctuations taken at fixed $\varphi_1$; the subscript $c$ signifies that all
diagrams disconnected from the insertion point $x$ must be disregarded (they are
already accounted for by the factor $\exp(-\mathcal A_1)$). Although these
expectation values depend on the slow modes $\varphi_1$, this dependence is
``quasilocal'', i.e.\ they can be written as series in powers of $\varphi_1(x)$
and its derivatives taken at the point $x$. Therefore the result of
$\tilde\varphi$ integration has the form
\[
\frac1{Z_1}\sum_\beta y_\alpha^\beta(L)\,\mathcal O_\beta\bigl(\varphi_1(x),\ \text{derivatives}\bigr)\,\varphi_1(x_1)\cdots\varphi_1(x_n)\,e^{-\mathcal A_1[\varphi_1]}\,,
\]
where in general the sum involves infinitely many terms, and $y_\alpha^\beta(L)$
are complicated sums of Feynman diagrams. Thus, in the step 1 of RG procedure the
insertion of the composite field $\mathcal O_\alpha(x)$ gets replaced by certain
linear combinations of slow-mode composite fields, with $L$-dependent
coefficients.

This structure holds true after we perform the step 2, which amounts to the change
of variables
\[
\varphi_1(x)=Z^{-\frac12}(L)\,\varphi_0(x/L)\,.
\]
As $\mathcal O_\beta$ in \eqref{eq:Ostep1} are polynomials of $\varphi_1$ and its
derivatives taken at $x$, we have
\begin{multline*}
\sum_\beta y_\alpha^\beta(L)\,\mathcal O_\beta\left(Z^{-\frac12}(L)\,\varphi_0(x/L),\ Z^{-\frac12}(L)\,\frac\partial{\partial x^\mu}\varphi_0(x/L),\dots\right)\\
=\sum_\beta z_\alpha{}^\beta(L)\,\mathcal O_\beta\left(\varphi_0(x/L),\ \frac\partial{\partial(x^\mu/L)}\varphi_0(x/L),\dots\right)\,,
\end{multline*}
i.e.\ under the full RG transformation
\begin{equation}\label{eq:RGfields}
\bigl\langle\mathcal O_\alpha(x)\cdots\bigr\rangle\big|_{\mathcal A}=\Bigl\langle\sum_\beta z_\alpha{}^\beta(L)\,\mathcal O_\beta(x/L)\cdots\Bigr\rangle\Big|_{\mathcal A'}\,,
\end{equation}
where $\mathcal A'$ is RG transformed action,
\[
\mathcal A'=RG_l\{\mathcal A\}\,,\qquad l=\log L\,.
\]
We will write
\begin{equation}\label{eq:RGop}
RG_{l,\mathcal A}\,\mathcal O_\alpha=\sum_\beta z_\alpha{}^\beta(L)\,\mathcal O_\beta\,,\qquad L=e^l\,,
\end{equation}
stressing the fact that $RG$ transformation acts on the fields linearly;
additional index $\mathcal A$ is written down to remind that this transformation
depends on the action $\mathcal A$.

Let us assume now that $\mathcal A$ is the fixed-point action,
\[
\mathcal A=\mathcal A_*\,.
\]
Then $\mathcal A'$ coincides with $\mathcal A$, and \eqref{eq:RGop} yields linear
relation for the correlation functions of the composite fields. Recall that by the
composition law of the RG transformations,
\[
RG_{l+l'}\{*\}=RG_l\{RG_{l'}\{*\}\}\,,
\]
in this case the coefficients $z$ must obey the equation
\[
z_\alpha{}^\beta(l+l')=\sum_\gamma z_\alpha{}^\gamma(l')\,z_\gamma{}^\beta(l)\,,
\]
where I changed the notations $z(l)=z(L=e^l)$. In other words these coefficients,
as the functions of $l$, must satisfy linear differential equation with constant
coefficients,
\[
\frac d{dl}\,z_\alpha{}^\beta(l)=-\sum_\gamma D_\alpha{}^\gamma\,z_\gamma{}^\beta(l)\,,
\qquad\text{where}\qquad
D_\alpha{}^\beta=-\frac d{dl}\,z_\alpha{}^\beta(l)\Big|_{l=0}\,.
\]
The coefficients $D_\alpha{}^\beta$ are matrix elements of the linear operator $D$
corresponding to the infinitesimal transformation \eqref{eq:RGop}
\[
RG_{\delta l,\mathcal A_*}=I-\delta l\,D+O(\delta l^2)\,.
\]
Let us denote $\Phi_\alpha(x)$ the eigenvectors of the operator $D$, i.e.\ special
linear combinations of the composite fields $\mathcal O_\alpha(x)$ which satisfy
the equations
\[
D\,\Phi_\alpha\equiv D_\alpha{}^\beta\,\Phi_\beta=D_\alpha\,\Phi_\alpha\,,
\]
where I use the notation $D_\alpha$ for the corresponding eigenvalues of $D$. Then
\[
RG_{l,\mathcal A_*}\,\Phi_\alpha=e^{-lD_\alpha}\,\Phi_\alpha\,,
\]
and for the fixed-point theory the equation \eqref{eq:RGop} reads
\begin{equation}\label{eq:scaling}
\bigl\langle\Phi_{\alpha_1}(x_1)\cdots\Phi_{\alpha_n}(x_n)\bigr\rangle\big|_{\mathcal A_*}
=L^{-D_{\alpha_1}}\cdots L^{-D_{\alpha_n}}\,\bigl\langle\Phi_{\alpha_1}(x_1/L)\cdots\Phi_{\alpha_n}(x_n/L)\bigr\rangle\big|_{\mathcal A_*}\,,
\end{equation}
i.e.\ the above correlation functions are homogeneous functions of the coordinates
$x_1,\dots,x_n$, of the degree $-(D_{\alpha_1}+\dots+D_{\alpha_n})$. For the
two-point correlation functions, which depend only on the distance between the two
points, it follows from \eqref{eq:scaling}
\[
\bigl\langle\Phi_\alpha(x)\Phi_\beta(x')\bigr\rangle\big|_{\mathcal A_*}=\frac{G_{\alpha,\beta}}{|x-x'|^{D_\alpha+D_\beta}}\,,
\]
where $G_{\alpha,\beta}$ are constants. In general, the equation expresses the
\emph{scale invariance} of the theory described by the fixed-point action
$\mathcal A_*$. The eigenvalues $D_\alpha$ are called the \emph{scale dimensions}
(or \emph{anomalous scale dimensions}) of the fields $\Phi_\alpha$. If
$\mathcal A_*$ is not a free-field theory, there is no reason to expect them to
coincide with ``canonical dimensions''. I'd like to stress that these anomalous
dimensions as well as the basic ``scaling'' fields $\Phi_\alpha$ are defined
relative to given fixed point $\mathcal A_*$. We will see that in the theory of
continuous phase transitions these dimensions determine the critical exponents.

We have seen that the space $\mathcal F$ is identical to the tangent space
$T\Sigma_{\mathcal A_*}$ and therefore the above operator $D$ must be closely
related to the operator $K$ which is the linearization of the RG generator
$B\{\mathcal A\}$ near the fixed point,
\begin{equation}\label{eq:BK}
B\{\mathcal A_*+\delta\mathcal A\}=K\,\delta\mathcal A\,.
\end{equation}
The relation is very simple. Any eigenvector $\Psi_\alpha$ of $K$,
\[
K\,\Psi_\alpha=\kappa_\alpha\,\Psi_\alpha\,,
\]
is the integral
\[
\Psi_\alpha=\int d^dx\,\Phi_\alpha(x)\,.
\]
Therefore the eigenvalues are related as
\[
\kappa_\alpha=d-D_\alpha\,,
\]
where $d$ is the number of the space dimensions. The term $d$ appears because in
the step 2 of the RG transformation of the integral $\int d^dx\,\Phi(x)$ we have to
make additional transformation of the integration measure $d^dx$,
\[
\int d^dx\,\Phi(x/L)=L^d\int d^dx\,\Phi(x)\,.
\]

As we have seen, the roles played by positive and negative eigenvalues of $K$ in
the formation of the RG flow near $\mathcal A_*$ are very different. The
negative-eigenvalue subspace $\mathcal F_-$ corresponds to the critical surface,
whereas the positive-eigenvalue subspace $\mathcal F_+$ is the linear
approximation to the unstable manifold $U(\mathcal A_*)\subset\Sigma(\infty)$. The
common terminology is as follows. One refers to
\[
\begin{array}{llll}
\Phi_\alpha & \text{with}\quad D_\alpha>d & \text{(i.e.\ with }\kappa_\alpha<0) & \text{as the ``irrelevant'' field,}\\
\Phi_\alpha & \text{with}\quad D_\alpha=d & \text{(i.e.\ with }\kappa_\alpha=0) & \text{as the ``marginal'' field,}\\
\Phi_\alpha & \text{with}\quad D_\alpha<d & \text{(i.e.\ with }\kappa_\alpha>0) & \text{as the ``relevant'' field.}
\end{array}
\]
If marginal fields are absent, the dimensionality of $U(\mathcal A_*)$ equals the
number of relevant fields.

The situation is more complex if the marginal fields are present. Let me assume
that there is only one marginal field $\Phi_M$. By definition, the relevant
(irrelevant) fields correspond to deformations $\delta\mathcal A$ of the
fixed-point action which increase (decrease) with the RG flow in the linear
approximation in $\delta\mathcal A$. The deformation
\[
\delta\mathcal A_M=\lambda_M\int d^dx\,\Phi_M(x)
\]
corresponding to the marginal field does not change in the linear approximation,
but it may well increase or decrease with the RG flow if we go beyond the linear
approximation, say take into account the terms $\sim\delta\mathcal A_M^2$. If
$\delta\mathcal A_M$ increases with the RG flow for $\lambda_M$ small enough, the
deformation $\delta\mathcal A_M$ (and the field $\Phi_M$) is called
\emph{marginally relevant}; similarly, if it decreases, one refers to $\Phi_M$ as
\emph{marginally irrelevant} field.

In many important respects the marginally relevant (marginally irrelevant) field
can be treated on the same footing with the relevant (irrelevant) ones. Thus, if a
marginally relevant field is present, the unstable manifold $U(\mathcal A_*)$ is
generated by the subspace $\mathcal F'_+$ which includes this marginally relevant
field along with all the relevant fields,
\[
\mathcal F'_+=TU(\mathcal A_*)|_{\mathcal A_*}\,.
\]
We are going to have more precise statements after studying some examples.

\subsection*{Gaussian fixed point}

Let us consider the simplest fixed point for which all calculations can be
performed explicitly, the so called \emph{Gaussian fixed point}. It is described
by the massless free action
\[
\mathcal A_{*,G}=\int d^dx\,\frac12\,(\partial\varphi_0)^2
\]
It is straightforward to check that it is indeed a fixed point. If we split
$\varphi_0$ into the slow and the fast components, these modes completely
decouple,
\[
\mathcal A_{*,G}=\int d^dx\,\left(\frac12\,(\partial\varphi_1)^2+\frac12\,(\partial\tilde\varphi)^2\right)\,,
\]
and therefore the step 1 of the RG procedure does not change the action
$\mathcal A_{*,G}$ at all,
\[
\mathcal A_1=\int d^dx\,\frac12\,(\partial\varphi_1)^2
\]
The step 2 is also easy to perform. Setting
\[
\varphi_1(x)=Z^{-\frac12}(L)\,\varphi_0(x/L)
\]
we obtain
\[
\mathcal A'[\varphi_0(x)]=\mathcal A_1\bigl[Z^{-\frac12}(L)\,\varphi_0(x/L)\bigr]\,,
\]
\[
\mathcal A'=Z^{-1}(L)\int d^dx\,\frac12\left(\frac\partial{\partial x^\mu}\varphi_0(x/L)\right)^2=\frac{L^{d-2}}{Z(L)}\int d^dx\,\frac12\left(\frac\partial{\partial x^\mu}\varphi_0(x)\right)^2\,.
\]
where the change of the coordinates $x\to Lx$ is made in the last step. Under the
choice
\[
Z^{\frac12}(L)=L^{\frac{d-2}2}
\]
$\mathcal A'$ coincides with the original action $\mathcal A_{*,G}$, so it indeed
is a fixed point.

It is not too difficult to find the RG action on the composite fields in this
fixed-point theory. It will be sufficient for us to consider the fields
\[
\mathcal O_n=\frac1{n!}\,\varphi_0^n\,.
\]
The step 1 of the RG transformation yields for $\mathcal O_n$
\[
\frac1{n!}\,\varphi_1^n+\sum_{k=1}^n\frac1{k!}\,\langle\tilde\varphi^k\rangle_{\tilde\varphi}\,\frac1{(n-k)!}\,\varphi_1^{n-k}\,.
\]
In the case of infinitesimal transformation $L=1+\delta l$ we have
\[
\langle\tilde\varphi^2\rangle_{\tilde\varphi}=\int_{\Lambda(1-\delta l)<|k|<\Lambda}\frac{d^dk}{(2\pi)^d}\,\frac1{k^2}=N_d\,\Lambda^{d-2}\,\delta l\,,
\]
where
\[
N_d=\frac{\Omega_d}{(2\pi)^d}\,,\qquad\text{and}\qquad\Omega_d=\frac{2\pi^{\frac d2}}{\Gamma(\frac d2)}
\]
is the area of $S^{d-1}$. Therefore we have in this case
\[
\frac1{n!}\,\varphi_1^n+\frac1{(n-2)!}\,\varphi_1^{n-2}\,\frac{N_d\Lambda^{d-2}}2\,\delta l\,.
\]
As this field does not involve derivatives the step 2 amounts to multiplying
$\varphi_1$ by $Z^{-\frac12}(L)\sim1-\frac{d-2}2\,\delta l$. Therefore the operator
$D$ acts on $\mathcal O_n$ as
\begin{equation}\label{eq:DOn}
D\,\mathcal O_n=n\,\frac{d-2}2\,\mathcal O_n-\frac{N_d\Lambda^{d-2}}2\,\mathcal O_{n-2}\,.
\end{equation}
It is not difficult to diagonalize this operator.

\begin{exercise}
Find the linear combinations
\[
\Phi_n=\mathcal O_n+\sum_{k=1}^na_k\,\mathcal O_{n-k}
\]
such that
\[
D\,\Phi_n=D_n\,\Phi_n\,.
\]
Show that
\begin{equation}\label{eq:Phi4}
\Phi_4=\frac1{4!}\,\varphi_0^4-\frac14\,\langle\varphi_0^2\rangle\,\varphi_0^2+\frac18\,\langle\varphi_0^2\rangle^2\,,
\qquad\langle\varphi_0^2\rangle=\int_{|k|<\Lambda}\frac{d^dk}{(2\pi)^d}\,\frac1{k^2}=\frac{N_d\Lambda^{d-2}}{d-2}\,,
\end{equation}
and generally
\[
\Phi_n=\frac1{n!}\,{:}\varphi_0^n{:}\ :=\ \frac1{n!}\,\varphi_0^n-\frac1{2!\,(n-2)!}\,\langle\varphi_0^2\rangle\,\varphi_0^{n-2}-\text{further self-contractions.}
\]
Calculate $\langle\Phi_n(x)\Phi_{n'}(0)\rangle$.
\end{exercise}

The eigenvalues $D_n$ are found immediately, without computing exact form of the
eigenvectors $\Phi_n$. Indeed, the eq.~\eqref{eq:DOn} shows that the operator $D$
acts on the fields $\mathcal O_n$ as triangular matrix (i.e.\ it mixes
$\varphi_0^n$ with itself and lower powers of $\varphi_0$), and so its eigenvalues
coincide with its diagonal elements. We find
\[
D_n=n\,\frac{d-2}2\,.
\]
Note that $D_1=\frac{d-2}2$ is the canonical dimension of $\varphi_0$ in
$d$-dimensional space. The dimensions $D_n$ coincide with the canonical dimensions
of the composite fields $\varphi_0^n$, which is not surprising because we are
considering the free massless field theory $\mathcal A_{*,G}$.

Dimensionality of the space $\mathcal F_+$ (i.e.\ the space of relevant fields)
associated with the fixed point $\mathcal A_{*,G}$ depends on $d$. The field
$\Phi_2=\frac12{:}\varphi_0^2{:}$ is always relevant; its dimension
\[
D_2=d-2<d\,.
\]
It remains the only relevant field for all $d>6$. For $d<6$ the field
$\Phi_3=\frac1{3!}{:}\varphi_0^3{:}$ with the dimension
\[
D_3=\frac32\,(d-2)
\]
becomes relevant, too. For $d<4$ it is joined by the field
$\Phi_4=\frac1{4!}{:}\varphi_0^4{:}$. The situation is summarized in the table
\begin{center}
\begin{tabular}{cc}
$d$ & relevant fields\\[4pt]
$d>6$ & $\varphi_0^2$\\
$6>d>4$ & $\varphi_0^2,\ \varphi_0^3$\\
$4>d>10/3$ & $\varphi_0^2,\ \varphi_0^3,\ \varphi_0^4$\\
$10/3>d>3$ & $\varphi_0^2,\ \varphi_0^3,\ \varphi_0^4,\ \varphi_0^5$\\
$\cdots$ & \\
$\frac{2n}{n-2}>d>\frac{2(n+1)}{n-1}$ & $\varphi_0^2,\ \varphi_0^3,\dots,\varphi_0^n$
\end{tabular}
\end{center}
At $d=2$ all polynomials of $\varphi_0$ are relevant.

Note that the above classification coincides with our earlier classification of
the perturbatively renormalizable scalar field theories. We see that in generic
dimensionality $d$ the perturbatively renormalizable field theories correspond to
the unstable manifold $U(\mathcal A_{*,G})$ of the gaussian fixed point. Also, we
realize now that this classification of renormalizable field theories is relevant
only to the vicinity of the gaussian fixed point (after all the perturbative
expansion is developed around this point); if we manage to find more complex fixed
point, its own unstable manifold will give rise to a family of finite local field
theories which may be not too similar to the perturbatively defined renormalizable
field theories we considered before.

Thus we have a correspondence between the unstable manifold $U(\mathcal A_{*,G})$
associated with the Gaussian fixed point and perturbatively renormalizable scalar
field theories in different dimensions $d$. This correspondence is more subtle if
$d$ takes special values such that one of the fields $\Phi_n$ acquires the
dimension $D_n=d$, i.e.\ becomes marginal. This happens for instance at $d=6$ with
$\Phi_3\sim\varphi_0^3$ and at $d=4$ with $\Phi_4\sim\varphi_0^4$. The last case is
of most interest to us, and we take it as main example.

\subsection*{\texorpdfstring{$\varphi^4$}{phi**4} flow at \texorpdfstring{$d=4$}{d=4}}

As at $d=4$ the field $\Phi_4$ is marginal, we have
\begin{equation}\label{eq:marginal4}
RG_l\Bigl\{\mathcal A_{*,G}+\lambda\int d^4x\,\Phi_4(x)\Bigr\}=\mathcal A_{*,G}+\lambda\int d^4x\,\Phi_4(x)+O(\lambda^2)\,,
\end{equation}
i.e.\ in the order $\lambda$ the perturbed action does not flow at all. However,
nontrivial flow can be generated beyond the linear approximation, in the higher
orders in $\lambda$. Let us estimate this flow in the order $\lambda^2$.

To perform the step 1 of the RG transformation we have to integrate over the fast
component $\tilde\varphi$ in
\[
\varphi_0=\varphi_1+\tilde\varphi\,.
\]
Consider the action $\mathcal A_0$ of the form
\begin{equation}\label{eq:A0phi4}
\mathcal A_0=\mathcal A_{*,G}+\lambda\int d^4x\,\Phi_4(x)=\int d^4x\,\left(\frac12\,(\partial\varphi_0)^2+\frac\lambda{4!}\,\varphi_0^4-\frac{\lambda\,b}4\,\varphi_0^2+\Const\right)
\end{equation}
where
\[
b=\langle\varphi_0^2\rangle_{\mathcal A_{*,G}}=\int_{|k|<\Lambda}\frac{d^4k}{(2\pi)^4\,k^2}=\frac{\Lambda^2}{16\pi^2}
\]
(recall that $\Phi_4=\frac1{4!}\varphi_0^4-\frac14\langle\varphi_0^2\rangle_{*,G}\,\varphi_0^2+\Const$,
as we should have found in the Exercise above). We have
\begin{multline}\label{eq:A1phi4}
e^{-\mathcal A_1[\varphi_1]}\sim e^{-\mathcal A_0[\varphi_1]}\times\int[D\tilde\varphi]\\
\times\exp\left\{-\int d^4x\,\left[\frac12\,(\partial\tilde\varphi)^2+\lambda\left(\frac1{4!}\,\tilde\varphi^4-\frac b4\,\tilde\varphi^2+\frac16\,\varphi_1\tilde\varphi^3+\frac14\,\varphi_1^2\tilde\varphi^2+\frac16\,\varphi_1^3\tilde\varphi\right)\right]\right\}\,.
\end{multline}
As we discussed already in Section~16, integration over $\tilde\varphi$ yields
\[
\mathcal A_1=\mathcal A_0-\text{sum of all connected diagrams;}
\]
the diagrams being built from the vertices
\[
\dia[0.45]{\draw[line width=0.9pt] (-1.4,0) -- (1.4,0); \node[vtx] at (0,0) {};}\qquad
\dia[0.45]{\foreach \a in {45,135,225,315} \draw[line width=0.9pt] (0,0) -- (\a:1.4); \node[vtx] at (0,0) {};}\qquad
\dia[0.45]{\draw[prop,densely dashed] (-1.4,0) -- (0,0); \foreach \a in {45,0,-45} \draw[line width=0.9pt] (0,0) -- (\a:1.4); \node[vtx] at (0,0) {};}\qquad
\dia[0.45]{\foreach \a in {135,225} \draw[prop,densely dashed] (0,0) -- (\a:1.4); \foreach \a in {45,315} \draw[line width=0.9pt] (0,0) -- (\a:1.4); \node[vtx] at (0,0) {};}\qquad
\dia[0.45]{\foreach \a in {150,180,210} \draw[prop,densely dashed] (0,0) -- (\a:1.4); \draw[line width=0.9pt] (0,0) -- (1.6,0); \node[vtx] at (0,0) {};}\,,
\]
where the solid lines correspond to the fast field $\tilde\varphi$ (and thus all
solid lines must be contracted), and the dashed lines represent the field
$\varphi_1$ which plays the role of an external field (i.e.\ the diagrams involve
no contractions of the dashed lines).

It is not difficult to check that the equation \eqref{eq:marginal4} indeed holds
to the order $\lambda$. The constant $b$ in the $\varphi^2$ term in
\eqref{eq:A0phi4} is such that the contribution of the diagram
\[
\dia[0.45]{\draw[prop,densely dashed] (-2,0) -- (2,0); \draw[line width=0.9pt] (0,0.6) circle (0.6); \node[vtx] at (0,0) {};}\ \sim\ \varphi_1^2
\]
is compensated exactly by the transformation of the term
$\frac{\lambda b}4\varphi_1^2$ in the step 2 of the RG procedure.

In the order $\lambda^2$ many new terms in $\mathcal A_1$ are generated. For
example the diagram
\begin{equation}\label{eq:tree6}
\diac[0.45]{\coordinate (L) at (-1.2,0); \coordinate (R) at (1.2,0); \draw[line width=0.9pt] (L) -- (R);
  \foreach \a in {120,180,240} \draw[prop,densely dashed] (L) -- ++(\a:1.1);
  \foreach \a in {60,0,-60} \draw[prop,densely dashed] (R) -- ++(\a:1.1);
  \node[vtx] at (L) {}; \node[vtx] at (R) {};}
\end{equation}
gives rise to the terms $\sim\varphi_1^6$. All such fields are irrelevant. The
diagram contributing to the terms $\sim\varphi_1^4$ is
\begin{equation}\label{eq:bubble4}
\diac[0.45]{\coordinate (L) at (-1,0); \coordinate (R) at (1,0); \draw[line width=0.9pt] (0,0) ellipse (1 and 0.7);
  \foreach \a in {135,225} \draw[prop,densely dashed] (L) -- ++(\a:1.1);
  \foreach \a in {45,-45} \draw[prop,densely dashed] (R) -- ++(\a:1.1);
  \node[vtx] at (L) {}; \node[vtx] at (R) {};}
\end{equation}
This diagram admits Taylor expansion in the external momenta. The zeroth term of
this expansion brings in the term $\varphi_1^4$ (without derivatives) into
$\mathcal A_1$, while the higher terms give rise to the irrelevant contributions
like $\varphi_1^2(\partial\varphi_1)^2$. I am going to use approximation in which
all irrelevant terms in $\mathcal A$ are neglected. Intuitive reason for this
approximation is that irrelevant terms die out very fast with the RG flow. Within
this approximation the diagram \eqref{eq:tree6} can be neglected, and the diagram
\eqref{eq:bubble4} can be evaluated at zero external momenta only, yielding the
following contribution to $\mathcal A_1$
\[
\frac{\Delta\lambda}{4!}\int d^4x\,\varphi_1^4(x)\,,
\]
where
\begin{equation}\label{eq:Dlambda}
\frac{\Delta\lambda}{4!}=-\frac{\lambda^2}{16}\int_{\Lambda>|k|>\Lambda/L}\frac{d^4k}{(2\pi)^4}\,\frac1{k^2}\,\frac1{k^2}\,.
\end{equation}
The factor $\lambda^2/16$ in this formula appears as follows. The term in
\eqref{eq:A1phi4} responsible for the vertices in the diagram \eqref{eq:bubble4}
is $\frac\lambda4\varphi_1^2\tilde\varphi^2$, so each vertex comes with the factor
$-\lambda/4$. One has to expand the exponential to the second order to get the
diagram \eqref{eq:bubble4}, this brings in additional factor $1/2!$. Finally, a
factor of 2 comes from the combinatorics of contractions in
$\langle\tilde\varphi^2\tilde\varphi^2\rangle$. The momentum integral in
\eqref{eq:Dlambda} is easy to evaluate, and we obtain
\[
\Delta\lambda=-3\,\frac{\lambda^2}{(4\pi)^2}\,\log L\,.
\]
In this order in $\lambda$ the step 2 is trivial. For the marginal term
$\int d^4x\,\varphi_1^4(x)$ the factor $Z^{-2}(L)$ which comes from the
transformation
\[
\varphi_1(x)=Z^{-\frac12}(L)\,\varphi_0(x/L)
\]
is exactly compensated by the transformation of the integration measure $d^4x$
(recall that $Z^{\frac12}=L$ at $d=4$; it is easy to see that possible corrections
to this formula are $\sim\lambda^2$).

Thus, in our approximation the result of the RG transformation of the action
\eqref{eq:A0phi4} again has the form \eqref{eq:A0phi4} with $\lambda$ replaced by
\[
\lambda'=\lambda-3\,\frac{\lambda^2}{(4\pi)^2}\,\log L\,.
\]
Taking $L=1+\delta l$ we see that $\lambda$ \emph{decreases} with the RG flow,
according to the equation
\begin{equation}\label{eq:WilsonBeta}
-\frac d{dl}\,\lambda(l)=3\,\frac{\lambda(l)^2}{(4\pi)^2}\,,
\end{equation}
i.e.\ at $d=4$ the field $\Phi_4$ is \emph{marginally irrelevant}. Note that the
r.h.s.\ of \eqref{eq:WilsonBeta} coincides with the $\lambda^2$ term
\eqref{eq:beta1} of the Callan-Symanzik $\beta$ function of $\varphi^4$ theory
which we obtained in Section~15.
